% Detection vs Forecasting Comparison Subsection
\subsection{Detection Versus Forecasting}
\label{sec:detection-vs-forecasting}

Detection and forecasting serve different clinical purposes. Detection aims to identify seizures as they occur for timely intervention. Forecasting aims to predict seizures before onset to allow preventive action. These different objectives lead to distinct technical approaches and validation requirements.

\subsubsection{Clinical Objectives}

Detection systems alert caregivers when a seizure occurs. This enables intervention during or immediately after the seizure. Detection is most valuable for generalized tonic-clonic or focal to bilateral tonic-clonic seizures where caregiver response can prevent injury or complications.

Forecasting systems provide advance warning of increased seizure likelihood. This allows patients to modify activities, take rescue medication, or avoid risky situations. Forecasting could improve quality of life by reducing the uncertainty of when seizures might occur.

\subsubsection{Clinical Maturity}

Detection shows greater clinical maturity than forecasting. Two detection studies meet Phase 3 FPR benchmarks. Commercial devices exist for detection including the Embrace watch and NightWatch armband. No forecasting study meets a clear clinical benchmark. No device has regulatory approval for seizure forecasting.

Detection benefits from clearer clinical requirements. The ILAE has defined guidelines for detection device validation including FPR targets. Forecasting lacks equivalent standardized guidelines. This gap complicates assessment of forecasting readiness.

\subsubsection{Reporting Practices}

Detection and forecasting differ in metric reporting. Forecasting studies consistently report patient-level success rates. Detection studies rarely report individual patient outcomes. This reporting gap complicates assessment of which patients benefit from each approach.

Forecasting studies report forecasting-specific metrics including IoC and TiW. Detection studies emphasize sensitivity and FPR. These different metrics reflect different clinical priorities but complicate cross-domain comparison.

\textbf{Key finding.} Detection is closer to clinical translation than forecasting. Two detection approaches meet clinical FPR benchmarks. Forecasting shows consistent AUC around 0.75 but lacks clear clinical benchmarks and regulatory pathways.

\textbf{Limitation identified.} Direct comparison between detection and forecasting is challenging due to different metrics, validation approaches, and clinical objectives. The field lacks standardized frameworks for assessing relative clinical value.
